\section{A bipartite family with all zeros on the canonical line}\label{sec:ehrhart}
For a graph $G$ on $v$ vertices let $P_G$ be the convex hull of
$\pm(e_i-e_j)$ over its edges, in the lattice $\mathbb Z^v\cap\{\sum x_i=0\}$.
Write $E_{m,n}(k)=|kP_{K_{m,n}}\cap\mathbb Z^{m+n}|$ for its Ehrhart polynomial.
The degree is $d=m+n-1$. Higashitani, Kummer and Michałek asked whether all zeros lie
on $\Re z=-1/2$ and whether decreasing a largest part gives interlacing, for
complete multipartite graphs (Conjecture~4.10, \href{https://arxiv.org/abs/1612.07538}{arXiv:1612.07538}).
A recent public manuscript reports a counterexample with six parts of size two
(\href{https://itsdavis.com/solved-conjectures/results/hkm-ehrhart/}{Cinna Davis, proof and verification}).
It leaves the complete bipartite question intact. Kölbl gave a conditional
result for $K_{4,n}$ (\href{https://arxiv.org/abs/2404.02136}{arXiv:2404.02136}, Theorem~3.3).
Carptopus then proved the line property through smaller part five using positive Jacobi
matrices (\href{https://doi.org/10.5281/zenodo.22486582}{2026 preprint}).
These are direct antecedents of the construction below.

\paragraph{Finite-width theorem.}
For $2\le m\le 32$ and every integer $n\ge m$, all the zeros of $E_{m,n}$
are simple and lie on $\Re z=-1/2$. The ordinates of the zeros strictly alternate with
those for $E_{m,n-1}$. Width $32$ is the last width for which we certify a positive Bezoutian.
The proof uses no zero-location theorem for analytic $L$-functions.

\subsection*{Reduction to two short polynomials}
Put $r=m-1$, $d=n+r$, and let $C_k$ denote the Ehrhart polynomial of a $k$-dimensional
cross-polytope. Its Ehrhart series is $(1+t)^k/(1-t)^{k+1}$.
Set $Q_k(y)=k!\,i^{-k}C_k((iy-1)/2)$. The cross-polytope recurrence
gives $Q_0=1$, $Q_1=y$, $Q_{k+1}=yQ_k-k^2Q_{k-1}$.
In particular, $Q_k$ is the characteristic polynomial of a real symmetric
tridiagonal matrix with consecutive off-diagonal entries $1,2,\ldots,k-1$.
The formula of Higashitani, Jochemko and Michałek for the $h^*$-polynomial
(\href{https://arxiv.org/abs/1807.07678}{arXiv:1807.07678}, Theorem~4.1)
has coefficients $\gamma_j=\binom{2j}{j}\binom{r}{j}\binom{n-1}{j}$ in the
basis $t^j(1+t)^{d-2j}$. Use $4t=(1+t)^2-(1-t)^2$ in the Ehrhart series.
Binomial expansion expresses $4^rE_{m,n}$ as
$\sum_{h=0}^r(-1)^h c_h C_{d-2h}$, where
$c_h=\sum_{j=h}^r4^{r-j}\gamma_j\binom{j}{h}$.
Thus the rotated polynomial has the exact expression
\[
 R_n(y)=4^r d!\,i^{-d} E_{m,n}((iy-1)/2)
       =\sum_{h=0}^r c_h(d)_{2h}Q_{d-2h}(y)
       =A_n(y)Q_n(y)-B_n(y)Q_{n-1}(y),
\]
where $(d)_k=d(d-1)\cdots(d-k+1)$. Expand each $Q_{d-2h}$ from
$Q_n,Q_{n-1}$ by forward or backward recurrence to define $A_n,B_n$.
Their $y$-degrees are at most $r,r-1$. $A_n$ has degree $r$ with
positive leading coefficient $c_0$. Their parities are $r,r-1$ respectively.
The construction also works at $n=m-1$, including a possible $Q_0$.

\subsection*{The positive form and the interlacing determinant}
Let $A=A_n$, $B=B_n$. Form the real symmetric matrix $H$ of coefficients in
$u^iv^j$ ($0\le i,j<r$) of the Bezoutian
\[
 \frac{A(u)B(v)-B(u)A(v)}{u-v}.
\]
Suppose $H$ is positive definite. Substitution $v=\bar u$ shows that
$A,B$ have no common zero and that $\Im(B(u)/A(u))<0$ for $\Im u>0$
where the quotient is defined. In fact $A(u)$ cannot vanish there, since
the diagonal Hermitian form at $(1,u,\ldots,u^{r-1})$ would vanish.
On the upper half-plane $Q_n/Q_{n-1}$ has positive imaginary part by
induction using $Q_{k+1}/Q_k=y-k^2/(Q_k/Q_{k-1})$.
The two ratios cannot coincide, so $R_n$ has only real zeros.
On the real axis the diagonal limit gives $W=A'B-AB'>0$.
Also $W_n=Q_n'Q_{n-1}-Q_nQ_{n-1}'>0$, since $W_1=1$ and
$W_n=Q_{n-1}^2+(n-1)^2W_{n-1}$ for $n\ge2$.

Set $C=y A_{n+1}-B_{n+1}$ and $D=n^2 A_{n+1}$, so that
$R_{n+1}=CQ_n-DQ_{n-1}$. The polynomial
\[
 T_n(y)=AD-BC=n^2 A_n A_{n+1}-y B_n A_{n+1}+B_n B_{n+1}
\]
is even in $y$. At every real root $t$ of $R_n$ there is a nonzero
$\lambda$ with $(Q_n,Q_{n-1})(t)=\lambda(B,A)(t)$.
Both vectors are nonzero by the recurrence and positive definiteness.
Differentiation gives $R_n'(t)=(\lambda^2W(t)+W_n(t))/\lambda$,
which proves simplicity. Moreover
\[
 \frac{R_{n+1}(t)}{R_n'(t)}
 =-\frac{\lambda^2 T_n(t)}{\lambda^2 W(t)+W_n(t)} .
\]
When $T_n$ is positive on the real axis, signs force a root of
$R_{n+1}$ in each interior gap and one in each exterior gap, since
both leading coefficients are positive. The degree difference is one,
so this accounts for all its roots, with multiplicity one.

\subsection*{Exact certification for unbounded $n$}
Only short polynomials involving $r$ occur in the two positivity tests.
The following details explain how finite calculations prove the universal
quantifier over $n$. In performing the backward recurrence, multiply
both coordinates by $S(n)=\prod_{j=1}^{r-1}(n-j)^2$.
Backward steps divide by distinct squared linear factors of $S$.
In the forward direction at most $r$ steps increase the degree
in $n$ by at most $2r$. In the backward direction the coordinates
have denominators dividing $S$ and, after multiplication by $S$,
degrees at most $2r$ (each step is $[yq_{k+1}-q_{k+2}]/(n+k+1)^2$,
where indices are relative to $n$). Thus each cleared coordinate in the
mixture has degree at most $7r$ in $n$:
$2r$ allows for $S$, $2r$ for forward recurrence, $r$ for $c_h$,
and $2r$ for $(d)_{2h}$. At $7r+1$ consecutive integers beginning at
$m$, the checker verifies that the coordinates divided by $S$ agree with
polynomials of degrees at most $r$ for $A_n$ and $r+1$ for $B_n$, by taking
forward differences. Multiplying these interpolants by $S$ and applying the
polynomial identity principle proves the degree bounds identically.

The Bezoutian splits into even and odd blocks by parity. Each entry has
degree at most $2r+1$ in $n$. For the leading minor of order $j$ of
either block, the checker computes its exact integer values at
$n=m,m+1,\ldots,m+(2r+1)j$ by fraction-free elimination.
All forward differences at $m$, through its actual degree, are positive.
For every integer $s\ge0$, Newton's identity
$F(m+s)=\sum_h\binom{s}{h}\Delta^h F(m)$ now proves strict positivity.
Sylvester's criterion applies to both blocks. At $n=m-1$ the two
blocks are also tested directly.

Finally, write $T_n(y)=U(n,y^2)$. Its degree in $n$ is at most $2r+2$.
The checker obtains $\Delta_n^h U(m,z)$ from $2r+3$ rows. For each nonzero
polynomial in $z$ it verifies a positive constant term and equal numbers
of Sturm sign changes at $0$ and $+\infty$. Hence each is positive
on $[0,\infty)$, and so is $U(m+s,z)$ for integral $s\ge0$ by
Newton's identity (the zeroth row is positive). The same Sturm test at
$n=m-1$ completes the square boundary case.
The scripts \texttt{verify\_ehrhart.py} and \texttt{verify\_interlacing.py}
perform the integer calculations. Each accepts the first and last widths
as arguments. Run both with \texttt{2 32}. Width one follows at once from
the cross-polytope recurrence. A proof for all complete bipartite graphs
would require a uniform argument controlling the forms or another way to
localize the roots as $r$ grows.
